% =====================================================================
% Appendix A.3 — Why no permutation test of the location effect is valid
%
% Regenerated by: audit/scripts/appendix_A3_permutation.py  (exit 0)
% Source code reused: audit/scripts/audit_1_2_permutation.py (signature
%   extraction, leave-one-out scorer, multiset counting)
% Audit item: 1.2
%
% PACKAGES REQUIRED: booktabs, amsmath (for \text in the display)
% PREAMBLE CLASH RISK: none. amsmath is almost certainly already loaded.
% Labels prefixed app:perm / tab:perm.
% =====================================================================

\section{Why no permutation test of the location effect is valid}
\label{app:perm}

Section~4.4 states that no valid permutation test of the location effect exists
for this design. The arithmetic behind that statement is set out here. It is a
diagnosis of the experimental design, not a retraction of the classification
result, which stands on the effect size given at the end of this section.

\subsection*{The two candidate permutation units}

A permutation test needs a unit that is exchangeable under the null hypothesis.
Two candidates are available, and each fails for a different reason.

The \emph{run-level} unit shuffles the location labels of individual severe runs.
The \emph{group-level} unit keeps the three replicates of a location together and
shuffles only the four location names.

Both spaces are small enough to enumerate exactly, so no sampling is needed
anywhere in this section.

\begin{table}[htbp]
\centering
\caption{Permutation spaces for the two fixed signature spaces. The run-level
count is the number of distinct labellings of the label multiset. The
partition-preserving count is the number of relabellings that map the partition
of runs into classes onto itself, which for a perfect classifier is exactly the
set of labellings reproducing the observed score.}
\label{tab:perm-spaces}
\begin{tabular}{lrrrrr}
\toprule
Signature space & Runs & Classes & Run-level space & Group-level space & Partition-preserving \\
\midrule
Two-mode $(f_1,f_3)$   & 11 & 4 & 92\,400 & 24 & 6 \\
Three-mode             &  9 & 3 &  1\,680 &  6 & 6 \\
\bottomrule
\end{tabular}
\end{table}

The run-level counts are multiset permutation counts. With class sizes
$3,3,3,2$ in the two-mode space and $3,3,3$ in the three-mode space,
%
\begin{equation*}
  \frac{11!}{3!\,3!\,3!\,2!} = 92\,400,
  \qquad
  \frac{9!}{3!\,3!\,3!} = 1\,680 .
\end{equation*}

\subsection*{Exact run-level $p$-values}

Both classifiers are perfect on the observed labelling, 11 of 11 and 9 of 9.
Because leave-one-out nearest-class-mean depends on the data only through the
partition of runs into classes, the labellings that reproduce a perfect score are
exactly those that preserve that partition, which are the relabellings permuting
equal-sized classes among themselves. There are $3!=6$ of them in each space.
The exact $p$-values are therefore closed-form rather than estimated:
%
\begin{equation*}
  p_{\text{two-mode}} = \frac{6}{92\,400} = 6.5\times10^{-5},
  \qquad
  p_{\text{three-mode}} = \frac{6}{1\,680} = 0.0036 .
\end{equation*}

\subsection*{Why neither unit gives a usable test}

\paragraph{Run-level permutation is anticonservative.} The three replicates of a
cell share one teardown, one damage application and one rebuild. They are
therefore correlated even under the null hypothesis that location carries no
information, and they are not exchangeable with runs drawn from other cells.
Treating them as independent narrows the null distribution and makes the
resulting $p$-value smaller than it should be. A defect of this kind can only
turn a true null into a false positive, never the reverse.

\paragraph{Group-level permutation has zero power by construction.} This is the
stronger statement and it does not depend on the data. The scorer is invariant
under the null's own relabelling operation: a group-level permutation permutes
class \emph{names} while leaving the partition of runs into classes identical, so
every distance the classifier computes is unchanged and the score cannot move.
Enumerating both spaces exhaustively confirms it. All 24 group-level labellings
in the two-mode space score 11 of 11, and all 6 in the three-mode space score 9
of 9. The attainable $p$-value is $24/24 = 6/6 = 1.0000$ in every case.

This must not be read as a null result. A null result would mean an effect was
tested for and not found. Here the test statistic cannot change no matter what
the data are, so the test would return the same value on a dataset with a real
location effect and on one with none. It carries no information about this result
or any other.

\paragraph{Consequently the design admits no valid test.} With three replicates
nested inside a single rebuild per cell, location cannot be separated from
rebuild. The only unit that is genuinely exchangeable under the null is the one
whose relabelling the statistic ignores.

\subsection*{The design change that would fix it}

Replicating the \emph{rebuild} rather than the taps. Three independent teardowns
and rebuilds per location, each with its own damage application, would make the
rebuild the exchangeable unit and give a permutation test with real power against
the location effect. This is a change to the experiment, not to the analysis, and
it is recorded as future work in Section~5.

\subsection*{The effect size reported instead}

With the $p$-values withdrawn, the separation of the class means carries the
result. Both signature spaces are reported because they measure different things:
the two-mode space includes Floor~3, whose second mode is voided by a harmonic,
and the three-mode space does not.

\begin{table}[htbp]
\centering
\caption{Separation of the closest pair of class means against the largest
distance from any run to its own class mean, in units of the mode-specific
$2\sigma$ reassembly floor.}
\label{tab:perm-effect}
\begin{tabular}{lrrr}
\toprule
Signature space & Closest class-mean pair & Largest within-class distance & Ratio \\
\midrule
Two-mode $(f_1,f_3)$, 11 runs, 4 classes & 26.6 & 2.06 & 12.9 \\
Three-mode, 9 runs, 3 classes            & 33.0 & 3.09 & 10.7 \\
\bottomrule
\end{tabular}
\end{table}

The closest pair is the base plate and Floor~1 in both spaces, which is the
expected result: they are physically adjacent and their plates soften overlapping
storeys. Every replicated severe run is assigned to its own location in both
spaces, and the closest confusable pair of classes is separated by more than ten
times the worst replicate spread.
